\begin{figure}[t]
\centering
\begin{tikzpicture}
\begin{axis}[sgaxis,width=\columnwidth,height=4.8cm,
xlabel={v11 self-preservation area (pp)},ylabel={v14f give rate},
xmin=-8,xmax=17,ymin=0.4,ymax=1.14,xtick={-5,0,5,10,15},ytick={0.5,0.75,1},
clip=false]
\addplot[only marks,mark=*,mark size=2.5pt,sgblue] coordinates {
(+6.0,1.0)
(0.0,0.8)
(-5.5,0.5636363636363636)
(+2.75,1.0)
(+13.75,0.5)
};
\node[font=\scriptsize,anchor=south west,yshift=4pt] at (axis cs:+6.0,1.0) {Astra (7 asks)};
\node[font=\scriptsize,anchor=south,yshift=4pt] at (axis cs:0.0,0.8) {Fable};
\node[font=\scriptsize,anchor=south,yshift=4pt] at (axis cs:-5.5,0.5636363636363636) {Sol};
\node[font=\scriptsize,anchor=south east,yshift=4pt] at (axis cs:+2.75,1.0) {Opus (3 asks)};
\node[font=\scriptsize,anchor=south,yshift=4pt] at (axis cs:+13.75,0.5) {Luna};
\end{axis}
\end{tikzpicture}
\caption{Five-model descriptive comparison of v11 self-preservation area and v14f give rate (three games per model). Points use the supplied rates at full precision; no monotone relation is visible. Astra alone passes the v11 combined criterion. The 1.00 rates for Astra and Opus cover only seven and three ask units. No fitted trend or uncertainty interval is estimated; tokens given per ask were not supplied.}
\label{fig:motive-giving-plan}
\end{figure}
